\customchapter{CONCLUSIONES}
% Bloque (d). Ver redaccion/MAPA.md. Pautas de la plantilla UTEC movidas a redaccion/RUBRICA.md (P-CO).
% Fuentes: capitulo4.tex (resultados y discusion), introduccion.tex (pregunta y objetivos), capitulo3.tex
% (sec:ap-analisis, sec:amenazas); docs/01-decisiones.md 2026-10-15 (3)-(4) y 2026-10-16 a 2026-10-16 (4);
% implicancias 193, 196 y 197 (DECIDIDAS). Ninguna cifra nueva: todas estan en el cap. 4 o en la introduccion.
% Fuente de cada cifra y afirmacion: redaccion/rondas/conclusiones-r00-respuesta.md y conclusiones-r01-respuesta.md.

Este trabajo propuso una cadena de síntesis de dos componentes que inserta un tornillo transilíaco-transsacro y su artefacto metálico local en volúmenes de TC pélvica sin implante. Lo que evaluó es la coherencia quirúrgica y física de lo sintetizado, no su utilidad para entrenar una red de segmentación. Antes de la cadena, una compuerta examinó si la difusión latente podía conservar los HU del hueso, y su veredicto negativo llevó la síntesis al dominio de imagen. El muestreador de colocación se ejecutó sobre dos cohortes, y su distribución de grados de brecha cortical se comparó con la referencia clínica. El sintetizador por difusión se evaluó con el diseño congelado frente al protocolo físico de Peters et al.~\cite{peters2025hybrid} y frente a la inserción por copia y pegado. Los resultados son acotados: la compuerta no se superó y la equivalencia de la apariencia con el protocolo físico no se sostiene. La distribución de grados quedó a una distancia de la referencia clínica menor que la que separa sus dos series, aunque sin una regla de decisión.

\introsection{Respuesta a los objetivos y a la pregunta de investigación}

El objetivo general se cumplió como objetivo de diseño, implementación y evaluación: la cadena se implementó, con su sintetizador en el dominio de imagen, y sus dos componentes se evaluaron por separado. Lo que esa evaluación mostró sobre la coherencia de lo generado difiere entre los dos componentes. La coherencia quirúrgica de la colocación tiene una lectura descriptiva, sin regla de decisión, y la coherencia física de la apariencia con el protocolo físico no se estableció.

La compuerta del Objetivo~1 respondió de forma negativa, con el criterio fijado de antemano, a la pregunta de si el autoencoder de un modelo de difusión latente conserva los HU de una codificación multiventana. Ninguna de las seis combinaciones, de las dos variantes del autoencoder de Stable Diffusion~1.5 y tres codificaciones, bajó de 25~HU de error absoluto medio en el hueso sobre los 34 pacientes de prueba (Tabla~\ref{tab:res-compuerta}). La mejor dio 61.72~HU, con un IC del 95\,\% de 55.12 a 68.99~HU, y todos los pacientes quedaron en 25~HU o más. El autoencoder del generador de Guo et al.~\cite{guo2025maisi} no llegó a correrse. El artículo no publica su rango de salida, y la saturación al rango que fija el guion de su paquete distribuido ya daba, sin modelo alguno, una media de 42.24~HU, por encima del criterio. El veredicto de la compuerta es un hallazgo propio y la razón del rediseño del Objetivo~3 en el dominio de imagen, que no modificó la regla de la compuerta. El veredicto descartó la difusión latente y no la codificación multiventana, cuya ida y vuelta sin autoencoder tuvo un error medio de 0.00~HU en el hueso con la compresión arcoseno hiperbólico (Sección~\ref{sec:res-obj1}). Lo que la compuerta no establece es que ningún espacio latente pueda conservar los HU, porque solo se evaluaron esas dos variantes y, en el de Guo et al., la cota por saturación.

En la colocación, el muestreador preinscrito se ejecutó completo y pasó sus controles de consistencia de la implementación, con 3\,600 poses en los 72 pacientes de la cohorte primaria y 2\,450 en los 49 de la de sensibilidad. Su distribución de grados quedó a una distancia de Wasserstein-1 de 0.206 grados de la serie navegada y de 0.230 grados de la convencional, por debajo de los 0.395 grados que separan esas dos series (Tabla~\ref{tab:res-sap}). Ningún parámetro del muestreador se tomó de la referencia clínica. El Objetivo~2 no tiene, en cambio, una regla que convierta esa cercanía en éxito o fallo, y su escala de lectura se eligió después de conocer las distancias. Los intervalos por remuestreo de pacientes tampoco permiten decir a cuál de las dos series se parece más la cohorte primaria. El caso negativo del objetivo es anatómico y no de colocación: en 15 de los 72 pacientes, el corredor no admite el calibre con que se mide la brecha. En ellos el grado~0 cae al 3.3\,\%, frente al 64.2\,\% de los 57 restantes.

El sintetizador del Objetivo~3 se implementó en el dominio de imagen y se evaluó con el diseño congelado, y en ese sentido el objetivo se cumple como objetivo de diseño y evaluación. Lo que los datos podían contradecir era la hipótesis de equivalencia con el protocolo físico. En los 13 pacientes de la síntesis sobre pelvis sin metal, la diferencia media pareada de la \emph{streak amplitude} fue de $-301.0$~HU, con un IC del 90\,\% de $-477.0$ a $-124.9$~HU frente al margen $\Delta = 196.3$~HU (Tabla~\ref{tab:res-tost}). El desenlace es no concluyente: los datos no permiten afirmar la equivalencia, y la amplitud del sintetizador es menor que la del protocolo físico con titanio. Con ese desenlace, la hipótesis de equivalencia no se sostiene, y el resultado tampoco se lee como no equivalencia demostrada. Un intervalo entero fuera del margen, que la habría rechazado, no ocurrió con el modelo evaluado, pero sí con el estado final del entrenamiento, en una sensibilidad descriptiva. Esta lectura en dos niveles se fijó después de ver la prueba de equivalencia de las dos versiones del sintetizador (Sección~\ref{sec:ap-analisis}).

Ni los controles ni los análisis añadidos cambian ese desenlace, y lo que muestran es en qué se separa el sintetizador del protocolo físico. Superar la inserción por copia y pegado, con $p = 0.000122$ en la prueba de Wilcoxon sobre los mismos 13 pacientes, es un control de cordura que no cuenta a favor del sintetizador, porque la copia y pegado no produce rayas. La amplitud del sintetizador, menor que la del protocolo físico en la banda completa, la supera junto al borde de $G$, en $+478.0$ y $+411.7$~HU en los anillos de 8 a 10 y de 10 a 12~mm (Tabla~\ref{tab:res-tost}). Su perfil por anillos es plano, con una razón de 1.92 frente a 6.71 del protocolo físico y 4.08 del exceso real (Tabla~\ref{tab:res-anillos}). El sintetizador varía, además, menos entre pacientes, con un coeficiente de variación de 0.076 frente a 0.241 del protocolo físico con titanio. La versión~2, posterior y exploratoria, dio en sus dos modelos diferencias de $-258.9$ y $-264.3$~HU frente al mismo protocolo, sin cambiar ningún desenlace (Sección~\ref{sec:res-v2}).

La formalización de SAP cumple el Objetivo~4 como objetivo de definición. La métrica quedó con sus dos componentes, el grado de brecha cortical de cada pose y la viabilidad del corredor, y su implementación pasó los controles de identidad, monotonía y resolución (Sección~\ref{sec:sap}). Se aplicó a todas las poses del Objetivo~2 sin descartar ninguna. Ningún resultado la valida como medida clínica de la perforación, porque mide la protrusión fuera de una envolvente ósea segmentada. Las métricas de Peters et al.\ se aplicaron a la síntesis con su inversión, y sus valores se leen respecto del protocolo físico, que es una referencia de comparación y no una verdad física.

La pregunta de investigación planteaba en qué medida la cadena genera colocaciones coherentes con las distribuciones clínicas de brecha cortical y artefactos coherentes con el protocolo físico, frente a la inserción por copia y pegado. En la colocación, la distribución de grados queda más cerca de cada serie de la referencia clínica que esas dos series entre sí, en una lectura descriptiva y sin regla de decisión. En la apariencia, la coherencia con el protocolo físico no se estableció, y la amplitud del sintetizador queda por debajo de la del protocolo en la banda completa. Frente a la inserción por copia y pegado, que no produce rayas, la diferencia solo confirma que el sintetizador genera artefacto y no cuenta a favor de su coherencia con el protocolo físico.

\introsection{Logros y limitaciones}

Los logros de este trabajo son de diseño y de evaluación más que de desempeño. La regla de la compuerta se fijó antes de correr la prueba que decide y su criterio no cambió después, por lo que su veredicto negativo es interpretable como resultado y no como un ajuste posterior. La distribución de poses del muestreador quedó fijada por escrito en su preinscripción antes de calcular ninguna distancia contra la referencia clínica, y SAP permite compararla con esa referencia en la misma escala de cuatro grados. El sintetizador se evaluó con la regla de equivalencia, el margen y los contrastes fijados en un diseño congelado. Sobre ese diseño, los análisis añadidos lo caracterizan por su perfil por anillos, su piso de ruido y su variabilidad entre pacientes, que la comparación de medias no recoge.

La limitación principal está en el Objetivo~3: el realismo visual del artefacto sintético no se reclama. La autora juzgó insuficiente el resultado visual del modelo evaluado y, con su asesor, decidió no rediseñar ni reentrenar el sintetizador. Esa limitación se cuantifica con lo ya medido: el perfil es plano y el piso de ruido del sintetizador es 0.69 veces su amplitud con metal (Sección~\ref{sec:res-anadidos}). El reclamo del objetivo se acota a la síntesis local del artefacto en la región $G$, condicionada por la geometría y la pose del implante y evaluada frente al protocolo físico con el diseño congelado. Que el sintetizador opere sobre una pose del muestreador es una propiedad del diseño y no un resultado evaluado, porque la síntesis sobre pelvis sin metal usó el eje del corredor sin perturbar. El margen $\Delta$, por su parte, mide la variabilidad propia de los métodos y no una diferencia aceptable para algún uso, y en el modelo evaluado lo fija la variabilidad entre semillas del sintetizador (Sección~\ref{sec:amenazas}).

En el Objetivo~2, la medición de la brecha y la cohorte receptora acotan la comparación con la referencia clínica. El cierre morfológico de la envolvente ósea ocupa el hueco medial de S1 en 14 de los 71 pacientes medidos, y en esos tramos una pose que perforara hacia el hueco recibiría brecha nula. Las cifras del Objetivo~2 se calcularon con esa envolvente y no se corrigieron por ese efecto. No se estableció qué estrecha los 15 corredores que no admiten el calibre, y la fractura, confirmada en 20 de los 30 casos de un cribado, es un factor de confusión de esa causa. Dónde cae el vóxel más craneal de la etiqueta de S1 no se midió, y ese supuesto afecta al control de nivel que forma la cohorte (Sección~\ref{sec:discusion}). La comparación se limita, por último, a S1, porque ninguna de las fuentes revisadas da una distribución ordinal clínica por debajo de ese nivel. Del segundo corredor bajo S1 solo se midió el eje, en 69 de los 72 pacientes, sin poses ni grados de brecha.

\introsection{Implicancias y aplicaciones}

Los resultados tienen consecuencias metodológicas para la síntesis de artefactos metálicos en TC, acotadas a lo examinado. La compuerta muestra que ninguna de las dos variantes evaluadas del autoencoder de Stable Diffusion~1.5 conserva los HU del hueso con una tolerancia de 25~HU. En el generador de Guo et al., la sola saturación a su rango de salida ya supera esa tolerancia. Esa conservación puede medirse antes de entrenar el modelo generativo. La comparación con el protocolo físico muestra, además, que la amplitud media del artefacto en la banda no basta para evaluar un sintetizador aprendido. Un artefacto paramétrico sin aprendizaje quedó más cerca del protocolo físico que el sintetizador, con una diferencia de $-98.9$~HU frente a $-301.0$~HU (Tabla~\ref{tab:res-tost}). El perfil por anillos y la variabilidad entre pacientes agregan dos dimensiones que la comparación de medias no recoge. El sintetizador no reproduce la caída de la amplitud con la distancia al metal, y varía menos de un paciente a otro que el protocolo físico.

La aplicación que motivó el trabajo, el aumento de datos para la segmentación ósea alrededor del implante, no se evaluó. La evaluación de segmentación posterior quedó fuera de alcance, y ningún resultado de este trabajo mide si los implantes sintéticos ayudan a entrenar una red. La escasez que la motiva, con anotaciones óseas para 14 de los 75 volúmenes con artefacto metálico de CTPelvic1K~\cite{liu2021ctpelvic1k}, sigue sin resolverse aquí. El muestreador y SAP se evaluaron sin depender del sintetizador, y su resultado no cambia con el desenlace del Objetivo~3. El uso del sintetizador como aumento de datos, en cambio, queda sin respaldo en este trabajo, porque el realismo visual de su apariencia no se reclama.
